\ifdefined\gkver\else\def\gkver{student}\fi
\documentclass[\gkver]{gaokaozhenti}
\begin{document}

\chapter{2020年山东}

\begin{enumerate}
\item
%% number: 1
%% paperName: 2020年山东等级考第 1 题 3 分
%% typeId: 1
%% type: 单选题
%% chapter: 运动和力
%% point: 超重失重
%% method: 无
%% score: 3
%% degree: 550
%% duplicateId: 0
%% body:
一质量为 $m$ 的乘客乘坐竖直电梯下楼，其位移 $s$ 与时间 $t$ 的关系图像如图所示。乘客所受支持力的大小用 $F_{N}$ 表示，速度大小用 $v$ 表示。重力加速度大小为 $g$。以下判断正确的是\xzanswer{D}

\onepicture[w=2.80cm, h=2.40cm]{figs/01.png}

\fourchoices[answer=D]
{$0\sim t_{1}$ 时间内，$v$ 增大，$F_{N}>mg$}
{$t_{1}\sim t_{2}$ 时间内，$v$ 减小，$F_{N}<mg$}
{$t_{2}\sim t_{3}$ 时间内，$v$ 增大，$F_{N}<mg$}
{$t_{2}\sim t_{3}$ 时间内，$v$ 减小，$F_{N}>mg$}

%% answer: D
%% memo:
\memoanswer{
A．由于 $s-t$ 图像的斜率表示速度，可知在 $0\sim t_{1}$ 时间内速度增加，即乘客的加速度向下，处于失重状态，则 $F_{N}<mg$，选项 A 错误；

B．在 $t_{1}\sim t_{2}$ 时间内速度不变，即乘客的匀速下降，则 $F_{N}=mg$，选项 B 错误；

CD．在 $t_{2}\sim t_{3}$ 时间内速度减小，即乘客的减速下降，超重，则 $F_{N}>mg$，选项 C 错误，D 正确；

故选 D。
}

\item
%% number: 2
%% paperName: 2020年山东等级考第 2 题 3 分
%% typeId: 1
%% type: 单选题
%% chapter: 电路
%% point: 电流强度
%% method: 无
%% score: 3
%% degree: 850
%% duplicateId: 0
%% body:
氚核 \ce{^{3}_{1}H} 发生 $\beta$ 衰变成为氦核 \ce{^{3}_{2}He}。假设含氚材料中 \ce{^{3}_{1}H} 发生 $\beta$ 衰变产生的电子可以全部定向移动，在 $3.2\times10^{4}\Us$ 时间内形成的平均电流为 $5.0\times10^{-8}\UA$。已知电子电荷量为 $1.6\times10^{-19}\UC$，在这段时间内发生 $\beta$ 衰变的氚核 \ce{^{3}_{1}H} 的个数为\xzanswer{B}

\fourchoices[answer=B]
{$5.0\times10^{14}$}
{$1.0\times10^{16}$}
{$2.0\times10^{16}$}
{$1.0\times10^{18}$}

%% answer: B
%% memo:
\memoanswer{
根据 $I=\frac{q}{t}=\frac{ne}{t}$ 可得产生的电子数为

$n=\frac{It}{e}=\frac{5.0\times10^{-8}\times3.2\times10^{4}}{1.6\times10^{-19}}=10^{16}$

个。

因在 $\beta$ 衰变中，一个氚核产生一个电子，可知氚核的个数为 $1.0\times10^{16}$ 个。

故选 B。
}

\item
%% number: 3
%% paperName: 2020年山东等级考第 3 题 3 分
%% typeId: 1
%% type: 单选题
%% chapter: 光学
%% point: 光的干涉
%% method: 无
%% score: 3
%% degree: 650
%% duplicateId: 0
%% body:
双缝干涉实验装置的截面图如图所示。光源 S 到 S$_{1}$、S$_{2}$ 的距离相等，O 点为 S$_{1}$、S$_{2}$ 连线中垂线与光屏的交点。光源 S 发出的波长为 $\lambda$ 的光，经 S$_{1}$ 出射后垂直穿过玻璃片传播到 O 点，经 S$_{2}$ 出射后直接传播到 O 点，由 S$_{1}$ 到 O 点与由 S$_{2}$ 到 O 点，光传播的时间差为 $\Delta t$。玻璃片厚度为 $10\lambda$，玻璃对该波长光的折射率为 1.5，空气中光速为 $c$，不计光在玻璃片内的反射。以下判断正确的是\xzanswer{A}

\onepicture[w=6.00cm, h=3.17cm]{figs/03.png}

\fourchoices[answer=A]
{$\Delta t=\frac{5\lambda}{c}$}
{$\Delta t=\frac{15\lambda}{2c}$}
{$\Delta t=\frac{10\lambda}{c}$}
{$\Delta t=\frac{15\lambda}{c}$}

%% answer: A
%% memo:
\memoanswer{
由几何知识知 $OS_{2}=OS_{1}$，设 $OS_{2}=x$，光从 S$_{2}$ 传播到 O 点的时间 $t_{1}=\frac{x}{c}$。由 $n=\frac{c}{v}$ 得，光在玻璃片中的传播速度 $v=\frac{c}{n}=\frac{c}{1.5}$，光垂直玻璃片入射，光的传播方向不变，则光从 S$_{1}$ 传播到 O 点的时间 $t_{2}=\frac{x-10\lambda}{c}+\frac{10\lambda}{v}=\frac{x+5\lambda}{c}$，所以 $\Delta t=t_{2}-t_{1}=\frac{5\lambda}{c}$。

故选 A。
}

\item
%% number: 4
%% paperName: 2020年山东等级考第 4 题 3 分
%% typeId: 1
%% type: 单选题
%% chapter: 机械波
%% point: 波的图象
%% method: 无
%% score: 3
%% degree: 550
%% duplicateId: 0
%% body:
一列简谐横波在均匀介质中沿 $x$ 轴负方向传播，已知 $x=\frac{5}{4}\lambda$ 处质点的振动方程为 $y=A\cos\left(\frac{2\pi}{T}t\right)$，则 $t=\frac{3}{4}T$ 时刻的波形图正确的是\xzanswer{D}

\fourchoices[answer=D, ispicture=true, h=2.00cm]
{figs/04a.png}
{figs/04b.png}
{figs/04c.png}
{figs/04d.png}

%% answer: D
%% memo:
\memoanswer{
$x=\frac{5}{4}\lambda$ 处质点的振动图像如图所示，$t=\frac{3}{4}T$ 时刻，$x=\frac{5}{4}\lambda$ 处质点的位移 $y=0$，质点处于平衡位置且向上运动，由于波沿 $x$ 轴负方向传播，可知符合要求的波形图为 D。

故选 D。

\onepicture[w=4.00cm, h=1.92cm]{figs/04_an.png}
}

\item
%% number: 5
%% paperName: 2020年山东等级考第 5 题 3 分
%% typeId: 1
%% type: 单选题
%% chapter: 交流电
%% point: 变压器
%% method: 无
%% score: 3
%% degree: 550
%% duplicateId: 0
%% body:
图甲中的理想变压器原、副线圈匝数比 $n_{1}:n_{2}=22:3$，输入端 a、b 所接电压 $u$ 随时间 $t$ 的变化关系如图乙所示。灯泡 L 的电阻恒为 $15\UO$，额定电压为 $24\UV$。定值电阻 $R_{1}=10\UO$、$R_{2}=5\UO$，滑动变阻器 $R$ 的最大阻值为 $10\UO$。为使灯泡正常工作，滑动变阻器接入电路的电阻应调节为\xzanswer{A}

\twopicture[numA=甲, numB=乙, wA=3.40cm, wB=6.50cm]{figs/05a.png}{figs/05b.png}

\fourchoices[answer=A]
{$1\UO$}
{$5\UO$}
{$6\UO$}
{$8\UO$}

%% answer: A
%% memo:
\memoanswer{
由题图乙可知，变压器原线圈两端电压的有效值为 $220\UV$。由理想变压器的电压关系 $\frac{U_{1}}{U_{2}}=\frac{n_{1}}{n_{2}}$ 得，变压器的输出电压 $U_{2}=\frac{n_{2}}{n_{1}}U_{1}=\frac{3}{22}\times220\UV=30\UV$。

灯泡正常工作，根据欧姆定律可知分压为 $24\UV$，则通过灯泡的电流，即副线圈部分的干路电流为 $I_{L}=\frac{U_{L}}{R_{L}}=\frac{24}{15}\UA=1.6\UA$。

根据串联分压规律可知，$R_{1}$ 和 $R_{2}$、$R$ 构成的并联电路部分的分压为 $U=U_{2}-U_{L}=30\UV-24\UV=6\UV$。

通过 $R_{1}$ 的电流为 $I_{1}=\frac{U}{R_{1}}=\frac{6}{10}\UA=0.6\UA$。

通过 $R_{2}$、$R$ 的电流为 $I_{2}=I_{L}-I_{1}=1.6\UA-0.6\UA=1\UA$。

$R_{2}$、$R$ 的分压为 $U=I_{2}(R_{2}+R)$，解得滑动变阻器的阻值为 $R=\frac{U}{I_{2}}-R_{2}=\frac{6}{1}\UO-5\UO=1\UO$。

A 正确，BCD 错误。

故选 A。
}

\item
%% number: 6
%% paperName: 2020年山东等级考第 6 题 3 分
%% typeId: 1
%% type: 单选题
%% chapter: 气体性质
%% point: 热力学第一定律
%% method: 无
%% score: 3
%% degree: 550
%% duplicateId: 0
%% body:
一定质量的理想气体从状态 a 开始，经 a$\to$b、b$\to$c、c$\to$a 三个过程后回到初始状态 a，其 $p-V$ 图像如图所示。已知三个状态的坐标分别为 a（$V_{0}$，$2p_{0}$）、b（$2V_{0}$，$p_{0}$）、c（$3V_{0}$，$2p_{0}$），以下判断正确的是\xzanswer{C}

\onepicture[w=3.40cm, h=2.72cm, align=right]{figs/06.png}

\fourchoices[answer=C]
{气体在 a$\to$b 过程中对外界做的功小于在 b$\to$c 过程中对外界做的功}
{气体在 a$\to$b 过程中从外界吸收的热量大于在 b$\to$c 过程中从外界吸收的热量}
{在 c$\to$a 过程中，外界对气体做的功小于气体向外界放出的热量}
{气体在 c$\to$a 过程中内能的减少量大于 b$\to$c 过程中内能的增加量}

%% answer: C
%% memo:
\memoanswer{
A．根据气体做功的表达式 $W=Fx=pSx=p\cdot\Delta V$ 可知，$p-V$ 图线和体积横轴围成的面积即为做功大小，所以气体在 a$\to$b 过程中对外界做的功等于 b$\to$c 过程中对外界做的功，A 错误；

B．气体从 a$\to$b，满足玻意耳定律 $pV=C$，所以 $T_{a}=T_{b}$，所以 $\Delta U_{ab}=0$，根据热力学第一定律 $\Delta U=Q+W$ 可知 $0=Q_{ab}+W_{ab}$；气体从 b$\to$c，温度升高，所以 $\Delta U_{bc}>0$，根据热力学第一定律可知 $\Delta U_{bc}=Q_{bc}+W_{bc}$，结合 A 选项可知 $W_{ab}=W_{bc}<0$，所以 $Q_{bc}>Q_{ab}$，b$\to$c 过程气体吸收的热量大于 a$\to$b 过程吸收的热量，B 错误；

C．气体从 c$\to$a，温度降低，所以 $\Delta U_{ca}<0$，气体体积减小，外界对气体做功，所以 $W_{ca}>0$，根据热力学第一定律可知 $Q_{ca}<0$，放出热量，C 正确；

D．理想气体的内能只与温度有关，根据 $T_{a}=T_{b}$ 可知 $\left|\Delta T_{ca}\right|=\left|\Delta T_{bc}\right|$，所以气体从 c$\to$a 过程中内能的减少量等于 b$\to$c 过程中内能的增加量，D 错误。

故选 C。
}

\item
%% number: 7
%% paperName: 2020年山东等级考第 7 题 3 分
%% typeId: 1
%% type: 单选题
%% chapter: 运动和力
%% point: 牛顿定律的应用
%% method: 无
%% score: 3
%% degree: 550
%% duplicateId: 0
%% body:
我国将在今年择机执行“天问 1 号”火星探测任务。质量为 $m$ 的着陆器在着陆火星前，会在火星表面附近经历一个时长为 $t_{0}$、速度由 $v_{0}$ 减速到零的过程。已知火星的质量约为地球的 0.1 倍，半径约为地球的 0.5 倍，地球表面的重力加速度大小为 $g$，忽略火星大气阻力。若该减速过程可视为一个竖直向下的匀减速直线运动，此过程中着陆器受到的制动力大小约为\xzanswer{B}

\fourchoices[answer=B]
{$m(0.4g-\frac{v_{0}}{t_{0}})$}
{$m(0.4g+\frac{v_{0}}{t_{0}})$}
{$m(0.2g-\frac{v_{0}}{t_{0}})$}
{$m(0.2g+\frac{v_{0}}{t_{0}})$}

%% answer: B
%% memo:
\memoanswer{
忽略星球的自转，万有引力等于重力 $G\frac{Mm}{R^{2}}=mg$，则 $\frac{g_{\text{火}}}{g_{\text{地}}}=\frac{M_{\text{火}}}{M_{\text{地}}}\cdot\frac{R_{\text{地}}^{2}}{R_{\text{火}}^{2}}=0.1\times\frac{1}{0.5^{2}}=0.4$，解得 $g_{\text{火}}=0.4g_{\text{地}}=0.4g$。

着陆器做匀减速直线运动，根据运动学公式可知 $0=v_{0}-at_{0}$，解得 $a=\frac{v_{0}}{t_{0}}$。

匀减速过程，根据牛顿第二定律得 $f-mg_{\text{火}}=ma$，解得着陆器受到的制动力大小为 $f=mg_{\text{火}}+ma=m\left(0.4g+\frac{v_{0}}{t_{0}}\right)$。

ACD 错误，B 正确。

故选 B。
}

\item
%% number: 8
%% paperName: 2020年山东等级考第 8 题 3 分
%% typeId: 1
%% type: 单选题
%% chapter: 力物体的平衡
%% point: 物体系平衡
%% method: 无
%% score: 3
%% degree: 450
%% duplicateId: 0
%% body:
如图所示，一轻质光滑定滑轮固定在倾斜木板上，质量分别为 $m$ 和 $2m$ 的物块 A、B，通过不可伸长的轻绳跨过滑轮连接，A、B 间的接触面和轻绳均与木板平行。A 与 B 间、B 与木板间的动摩擦因数均为 $\mu$，设最大静摩擦力等于滑动摩擦力。当木板与水平面的夹角为 $45^\circ$ 时，物块 A、B 刚好要滑动，则 $\mu$ 的值为\xzanswer{C}

\onepicture[w=3.60cm, h=3.22cm, align=right]{figs/08.png}

\fourchoices[answer=C]
{$\frac{1}{3}$}
{$\frac{1}{4}$}
{$\frac{1}{5}$}
{$\frac{1}{6}$}

%% answer: C
%% memo:
\memoanswer{
当木板与水平面的夹角为 $45^\circ$ 时，两物块刚好滑动，对 A 物块受力分析如图。沿斜面方向，A、B 之间的滑动摩擦力 $f_{1}=\mu N=\mu mg\cos45^\circ$，根据平衡条件可知 $T=mg\sin45^\circ+\mu mg\cos45^\circ$。

对 B 物块受力分析如图。沿斜面方向，B 与斜面之间的滑动摩擦力 $f_{2}=\mu N'=\mu\cdot3mg\cos45^\circ$，根据平衡条件可知 $2mg\sin45^\circ=T+\mu mg\cos45^\circ+\mu\cdot3mg\cos45^\circ$。

两式相加，可得 $2mg\sin45^\circ=mg\sin45^\circ+\mu mg\cos45^\circ+\mu mg\cos45^\circ+\mu\cdot3mg\cos45^\circ$，解得 $\mu=\frac{1}{5}$。

故选 C。

\twopicture[wA=2.20cm, wB=2.20cm, align=right]{figs/08_an1.png}{figs/08_an2.png}
}

\item
%% number: 9
%% paperName: 2020年山东等级考第 9 题 4 分
%% typeId: 2
%% type: 多选题
%% chapter: 光学
%% point: 光的折射
%% method: 无
%% score: 4
%% degree: 650
%% duplicateId: 0
%% body:
截面为等腰直角三角形的三棱镜如图甲所示。DE 为嵌在三棱镜内部紧贴 BB'C'C 面的线状单色可见光光源，DE 与三棱镜的 ABC 面垂直，D 位于线段 BC 的中点。图乙为图甲中 ABC 面的正视图。三棱镜对该单色光的折射率为 $\sqrt{2}$，只考虑由 DE 直接射向侧面 AA'C'C 的光线。下列说法正确的是\xzanswer{AC}

\twopicture[numA=甲, numB=乙, wA=6.00cm, wB=4.60cm]{figs/09a.png}{figs/09b.png}

\fourchoices[answer=AC]
{光从 AA'C'C 面出射的区域占该侧面总面积的 $\frac{1}{2}$}
{光从 AA'C'C 面出射的区域占该侧面总面积的 $\frac{2}{3}$}
{若 DE 发出的单色光频率变小，AA'C'C 面有光出射的区域面积将增大}
{若 DE 发出的单色光频率变小，AA'C'C 面有光出射的区域面积将减小}

%% answer: AC
%% memo:
\memoanswer{
AB．由题可知 $\sin C=\frac{1}{\sqrt{2}}$，可知临界角为 $45^\circ$，因此从 D 点发出的光，竖直向上从 M 点射出的光线恰好是出射光线的边缘，同时 C 点也恰好是出射光线的边缘，如图所示，因此光线只能从 MC 段射出，根据几何关系可知，M 恰好为 AC 的中点，因此在 AA'C'C 平面上有一半的面积有光线射出，A 正确，B 错误；

C．由于频率越高，折射率越大，当光源发出的光的频率变小，折射率也会变小，导致临界角会增大，这时 M 点上方也会有光线出射，因此出射光线区域的面积将增大，C 正确，D 错误。

故选 AC。

\onepicture[w=4.00cm, h=3.41cm]{figs/09_an.png}
}

\item
%% number: 10
%% paperName: 2020年山东等级考第 10 题 4 分
%% typeId: 2
%% type: 多选题
%% chapter: 电场
%% point: 电势能电势差电场力做功
%% method: 无
%% score: 4
%% degree: 650
%% duplicateId: 0
%% body:
真空中有两个固定的带正电的点电荷，电荷量不相等。一个带负电的试探电荷置于二者连线上的 O 点时，仅在电场力的作用下恰好保持静止状态。过 O 点作两正电荷连线的垂线，以 O 点为圆心的圆与连线和垂线分别交于 a、c 和 b、d，如图所示。以下说法正确的是\xzanswer{BD}

\onepicture[w=5.00cm, h=2.70cm, align=right]{\includesvg{figs/10}}

\fourchoices[answer=BD]
{a 点电势低于 O 点}
{b 点电势低于 c 点}
{该试探电荷在 a 点的电势能大于在 b 点的电势能}
{该试探电荷在 c 点的电势能小于在 d 点的电势能}

%% answer: BD
%% memo:
\memoanswer{
带负电的试探电荷仅在电场力的作用下恰好保持静止状态，由 $F=qE$ 知，O 点的电场强度为零，则两正点电荷连线上的电场线方向如图所示，由沿电场线方向电势逐渐降低知，a 点的电势高于 O 点，故 A 错误；

设左边正点电荷所带的电荷量为 $q_{1}$，右边正点电荷所带的电荷量为 $q_{2}$，则在 O 点有 $k\frac{q_{1}}{r_{1}^{2}}=k\frac{q_{2}}{r_{2}^{2}}$，在 b 点处平行两正电荷连线方向有 $k\frac{q_{1}}{\left(\frac{r_{1}}{\cos\theta_{1}}\right)^{2}}\cdot\cos\theta_{1}=\frac{kq_{1}\cos^{3}\theta_{1}}{r_{1}^{2}}>k\frac{q_{2}}{\left(\frac{r_{2}}{\cos\theta_{2}}\right)^{2}}\cos\theta_{2}=\frac{kq_{2}\cos^{3}\theta_{2}}{r_{2}^{2}}$，故 b 点电场强度平行两正电荷连线方向的分量指向右方，则 b 点场强方向斜向右上方，由沿电场线方向电势逐渐降低知，O 点电势低于 a、c 两点，高于 b、d 两点，故 b 点的电势低于 a、c 两点，由负电荷在电势低处电势能大知，试探电荷在 a 点的电势能小于在 b 点的电势能，故 B 正确 C 错误；

由对称性知，d 点的电势低于 c 点，由上述分析可知，带负电的试探电荷在 c 点的电势能小于在 d 点的电势能，故 D 正确。

故选 BD。

\onepicture[w=5.00cm, h=2.87cm]{figs/10_an.png}
}

\item
%% number: 11
%% paperName: 2020年山东等级考第 11 题 4 分
%% typeId: 2
%% type: 多选题
%% chapter: 运动和力
%% point: 变加速直线运动
%% method: 无
%% score: 4
%% degree: 450
%% duplicateId: 0
%% body:
如图所示，质量为 $M$ 的物块 A 放置在光滑水平桌面上，右侧连接一固定于墙面的水平轻绳，左侧通过一倾斜轻绳跨过光滑定滑轮与一竖直轻弹簧相连。现将质量为 $m$ 的钩码 B 挂于弹簧下端，当弹簧处于原长时，将 B 由静止释放，当 B 下降到最低点时（未着地），A 对水平桌面的压力刚好为零。轻绳不可伸长，弹簧始终在弹性限度内，物块 A 始终处于静止状态。以下判断正确的是\xzanswer{ACD}

\onepicture[w=3.40cm, h=3.30cm, align=right]{figs/11.png}

\fourchoices[answer=ACD]
{$M<2m$}
{$2m<M<3m$}
{在 B 从释放位置运动到最低点的过程中，所受合力对 B 先做正功后做负功}
{在 B 从释放位置运动到速度最大的过程中，B 克服弹簧弹力做的功等于 B 机械能的减少量}

%% answer: ACD
%% memo:
\memoanswer{
AB．由题意可知 B 物体可以在开始位置到最低点之间做简谐振动，故在最低点时有弹簧弹力 $T=2mg$；对 A 分析，设绳子与桌面间夹角为 $\theta$，则依题意有 $2mg\sin\theta=Mg$，故有 $M<2m$，故 A 正确，B 错误；

C．由题意可知 B 从释放位置到最低点过程中，开始弹簧弹力小于重力，物体加速，合力做正功；后来弹簧弹力大于重力，物体减速，合力做负功，故 C 正确；

D．对于 B，在从释放到速度最大过程中，B 机械能的减少量等于弹簧弹力所做的负功，即等于 B 克服弹簧弹力所做的功，故 D 正确。

故选 ACD。
}

\item
%% number: 12
%% paperName: 2020年山东等级考第 12 题 3 分
%% typeId: 2
%% type: 多选题
%% chapter: 电磁感应
%% point: 电磁感应受力分析
%% method: 无
%% score: 3
%% degree: 350
%% duplicateId: 0
%% body:
如图所示，平面直角坐标系的第一和第二象限分别存在磁感应强度大小相等、方向相反且垂直于坐标平面的匀强磁场，图中虚线方格为等大正方形。一位于 $Oxy$ 平面内的刚性导体框 abcde 在外力作用下以恒定速度沿 $y$ 轴正方向运动（不发生转动）。从图示位置开始计时，4 s 末 bc 边刚好进入磁场。在此过程中，导体框内感应电流的大小为 $I$，ab 边所受安培力的大小为 $F_{ab}$，二者与时间 $t$ 的关系图像，可能正确的是\xzanswer{BC}

\onepicture[w=5.00cm, h=4.98cm, align=right]{figs/12.png}

\fourchoices[answer=BC, ispicture=true, h=2.30cm]
{figs/12a.png}
{figs/12b.png}
{figs/12c.png}
{figs/12d.png}

%% answer: BC
%% memo:
\memoanswer{
AB．因为 4 s 末 bc 边刚好进入磁场，可知线框的速度每秒运动一个方格，故在 $0\sim1\Us$ 内只有 ae 边切割磁场，设方格边长为 $L$，根据 $E_{1}=2BLv$，$I_{1}=\frac{E_{1}}{R}$，可知电流恒定；2 s 末时线框在第二象限长度最长，此时有 $E_{2}=3BLv$，$I_{2}=\frac{E_{2}}{R}$，可知 $I_{2}=\frac{3}{2}I_{1}$；$2\sim4\Us$ 线框有一部分进入第一象限，电流减小，在 4 s 末同理可得 $I_{3}=\frac{1}{2}I_{1}$。综上分析可知 A 错误，B 正确；

CD．根据 $F_{ab}=BIL_{ab}$ 可知在 $0\sim1\Us$ 内 ab 边所受的安培力线性增加；1 s 末安培力为 $F_{ab}=BI_{1}L$，在 2 s 末可得安培力为 $F_{ab}'=B\times\frac{3}{2}I_{1}\times2L$，所以有 $F_{ab}'=3F_{ab}$；由图像可知 C 正确，D 错误。

故选 BC。
}

\item
%% number: 13
%% paperName: 2020年山东等级考第 13 题 6 分
%% typeId: 4
%% type: 实验题
%% chapter: 直线运动
%% point: 运动图象
%% method: 无
%% score: 6
%% degree: 550
%% duplicateId: 0
%% body:
2020 年 5 月，我国进行了珠穆朗玛峰的高度测量，其中一种方法是通过使用重力仪测量重力加速度，进而间接测量海拔高度。某同学受此启发就地取材设计了如下实验，测量当地重力加速度的大小。实验步骤如下：

\begin{steps}
\item
如图甲所示，选择合适高度的垫块，使木板的倾角为 $53^\circ$，在其上表面固定一与小物块下滑路径平行的刻度尺（图中未画出）。
\item
调整手机使其摄像头正对木板表面，开启视频录像功能。将小物块从木板顶端释放，用手机记录下小物块沿木板向下做加速直线运动的情况。然后通过录像的回放，选择小物块运动路径上合适的一点作为测量参考点，得到小物块相对于该点的运动距离 $L$ 与运动时间 $t$ 的数据。
\item
该同学选取部分实验数据，画出了 $\frac{2L}{t}-t$ 图像，利用图像数据得到小物块下滑的加速度大小为 $5.6\Umsq$。
\item
再次调节垫块，改变木板的倾角，重复实验。
\end{steps}

回答以下问题：
\begin{enumerate}
\item
当木板的倾角为 $37^\circ$ 时，所绘图像如图乙所示。由图像可得，物块过测量参考点时速度的大小为\rule{1.8cm}{0.4pt}$\Ums$；选取图线上位于坐标纸网格交叉点上的 A、B 两点，利用 A、B 两点数据得到小物块下滑加速度的大小为\rule{1.8cm}{0.4pt}$\Umsq$（结果均保留 2 位有效数字）。

\item
根据上述数据，进一步分析得到当地的重力加速度大小为\rule{1.8cm}{0.4pt}$\Umsq$（结果保留 2 位有效数字，$\sin37^\circ=0.60$，$\cos37^\circ=0.80$）。
\end{enumerate}

\twopicture[numA=甲, numB=乙, wA=6.00cm, wB=8.00cm]{figs/13a.png}{figs/13b.png}

%% answer:
\jdanswer{
\begin{enumerate}
\item
0.32 或 0.33，3.1

\item
9.4

\end{enumerate}
}
%% memo:
\memoanswer{
（1）由 $L=v_{0}t+\frac{1}{2}at^{2}$，得 $\frac{2L}{t}=2v_{0}+at$，则 $\frac{2L}{t}-t$ 图线的斜率表示加速度，截距表示 $2v_{0}$，由题图乙得 $2v_{0}=65\times10^{-2}\Ums$，解得物块过测量参考点的速度大小为 $v_{0}=0.33\Ums$，小物块下滑的加速度大小 $a=\frac{(185-70)\times10^{-2}\Ums}{(39.5-2.0)\times10^{-2}\Us}=3.1\Umsq$。

（2）当木板的倾角为 $37^\circ$ 时，由牛顿第二定律得 $mg\sin37^\circ-\mu mg\cos37^\circ=ma$；同理，当木板的倾角为 $53^\circ$ 时，有 $mg\sin53^\circ-\mu mg\cos53^\circ=ma_{1}$，由题知 $a_{1}=5.6\Umsq$，联立解得 $g=9.4\Umsq$。
}

\item
%% number: 14
%% paperName: 2020年山东等级考第 14 题 8 分
%% typeId: 4
%% type: 实验题
%% chapter: 电路
%% point: 测电动势与内阻
%% method: 无
%% score: 8
%% degree: 550
%% duplicateId: 0
%% body:
实验方案对实验测量的精度有直接的影响，某学习小组对“测量电源的电动势和内阻”的实验方案进行了探究。实验室提供的器材有：

干电池一节（电动势约 $1.5\UV$，内阻小于 $1\UO$）；

电压表 V（量程 $3\UV$，内阻约 $3\UkO$）；

电流表 A（量程 $0.6\UA$，内阻约 $1\UO$）；

滑动变阻器 $R$（最大阻值为 $20\UO$）；

定值电阻 $R_{1}$（阻值 $2\UO$）；

定值电阻 $R_{2}$（阻值 $5\UO$）；

开关一个，导线若干。
\begin{enumerate}
\item
该小组按照图甲所示的电路进行实验，通过调节滑动变阻器阻值使电流表示数逐渐接近满偏，记录此过程中电压表和电流表的示数，利用实验数据在 $U-I$ 坐标纸上描点，如图乙所示，结果发现电压表示数的变化范围比较小，出现该现象的主要原因是\rule{1.5cm}{0.4pt}。（单选，填正确答案标号）

A．电压表分流

B．干电池内阻较小

C．滑动变阻器最大阻值较小

D．电流表内阻较小

\twopicture[numA=甲, numB=乙, wA=3.60cm, wB=7.50cm]{figs/14a.png}{figs/14b.png}

\item
针对电压表示数的变化范围比较小的问题，该小组利用实验室提供的器材改进了实验方案，重新测量得到的数据如下表所示。

\begin{center}
\begin{tabular}{|c|c|c|c|c|c|c|c|}
\hline
序号 & 1 & 2 & 3 & 4 & 5 & 6 & 7 \\
\hline
$I/\UA$ & 0.08 & 0.14 & 0.20 & 0.26 & 0.32 & 0.36 & 0.40 \\
\hline
$U/\UV$ & 1.35 & 1.20 & 1.05 & 0.88 & 0.73 & 0.71 & 0.52 \\
\hline
\end{tabular}
\end{center}

请根据实验数据，回答以下问题：

①答题卡的坐标纸上已标出后 3 组数据对应的坐标点，请在答题卡的坐标纸上标出前 4 组数据对应的坐标点并画出 $U-I$ 图像\rule{1.5cm}{0.4pt}。

\onepicture[w=6.50cm, h=3.98cm]{figs/14c.png}

②根据实验数据可知，所选的定值电阻为\rule{1.5cm}{0.4pt}（填“$R_{1}$”或“$R_{2}$”）。

③用笔画线代替导线，请在答题卡上按照改进后的方案，将实物图连接成完整电路。

\onepicture[w=5.50cm, h=3.48cm]{figs/14d.png}
\end{enumerate}

%% answer:
\jdanswer{
\begin{enumerate}
\item
B

\item
①见解析图；②$R_{1}$；③见解析图

\end{enumerate}
}
%% memo:
\memoanswer{
（1）由闭合电路欧姆定律得，电压表的示数为 $U=\frac{E}{R+r}R=\frac{1}{1+\frac{r}{R}}E$。电压表的示数变化范围较小，可能是因为干电池内阻较小，也可能是滑动变阻器最大阻值较小，而根据题图乙中数据可知电压表示数一直比较接近电源电动势，则可排除滑动变阻器最大阻值较小的可能，故选 B。

（2）①根据表格中的数据描点连线，用一条直线将尽可能多的点连起来，不在直线上的点尽量对称分布在直线两侧，偏远的点舍去，图像如图所示。

\onepicture[w=7.00cm, h=4.32cm]{figs/14_an1.png}

②由 $U-I$ 图像知，干电池的电动势为 $E=1.55\UV$，电压表示数为 0 时，电路中的电流为 $0.6\UA$，由闭合电路欧姆定律得 $r+R=\frac{E}{I}=\frac{1.55\UV}{0.6\UA}\approx2.58\UO$，电源内阻小于 $1\UO$，故所选的定值电阻为 $R_{1}$。

③实物图连接如图所示。

\onepicture[w=5.50cm, h=3.25cm]{figs/14_an2.png}
}

\item
%% number: 15
%% paperName: 2020年山东等级考第 15 题 7 分
%% typeId: 6
%% type: 计算题
%% chapter: 气体性质
%% point: 理想气体状态方程
%% method: 无
%% score: 7
%% degree: 550
%% duplicateId: 0
%% body:
中医拔罐的物理原理是利用玻璃罐内外的气压差使罐吸附在人体穴位上，进而治疗某些疾病。常见拔罐有两种，如图所示，左侧为火罐，下端开口；右侧为抽气拔罐，下端开口，上端留有抽气阀门。使用火罐时，先加热罐中气体，然后迅速按到皮肤上，自然降温后火罐内部气压低于外部大气压，使火罐紧紧吸附在皮肤上。抽气拔罐是先把罐体按在皮肤上，再通过抽气降低罐内气体压强。某次使用火罐时，罐内气体初始压强与外部大气压相同，温度为 $450\UK$，最终降到 $300\UK$，因皮肤凸起，内部气体体积变为罐容积的 $\frac{20}{21}$。若换用抽气拔罐，抽气后罐内剩余气体体积变为抽气拔罐容积的 $\frac{20}{21}$，罐内气压与火罐降温后的内部气压相同。罐内气体均可视为理想气体，忽略抽气过程中气体温度的变化。求应抽出气体的质量与抽气前罐内气体质量的比值。

\twopicture[capA=火罐, capB=抽气拔罐, showprefix=false, wA=3.20cm, wB=3.60cm]{figs/15a.png}{figs/15b.png}

%% answer:
\jdanswer{
$\frac{\Delta m}{m}=\frac{1}{3}$
}
%% memo:
\memoanswer{
设火罐内气体初始状态参量分别为 $p_{1}$、$T_{1}$、$V_{1}$，温度降低后状态参量分别为 $p_{2}$、$T_{2}$、$V_{2}$，罐的容积为 $V_{0}$，大气压强为 $p_{0}$，由题意知 $p_{1}=p_{0}$，$T_{1}=450\UK$，$V_{1}=V_{0}$，$T_{2}=300\UK$，$V_{2}=\frac{20}{21}V_{0}$。

由理想气体状态方程得 $\frac{p_{0}V_{0}}{T_{1}}=\frac{p_{2}\cdot\frac{20}{21}V_{0}}{T_{2}}$，代入数据得 $p_{2}=0.7p_{0}$。

对于抽气罐，设初态气体状态参量分别为 $p_{3}$、$V_{3}$，末态气体状态参量分别为 $p_{4}$、$V_{4}$，罐的容积为 $V_{0}'$，由题意知 $p_{3}=p_{0}$，$V_{3}=V_{0}'$，$p_{4}=p_{2}$。

由玻意耳定律得 $p_{0}V_{0}'=p_{2}V_{4}$，联立并代入数据得 $V_{4}=\frac{10}{7}V_{0}'$。

设抽出的气体的体积为 $\Delta V$，由题意知 $\Delta V=V_{4}-\frac{20}{21}V_{0}'$，故应抽出气体的质量与抽气前罐内气体质量的比值为 $\frac{\Delta m}{m}=\frac{\Delta V}{V_{4}}$，联立并代入数据得 $\frac{\Delta m}{m}=\frac{1}{3}$。
}

\item
%% number: 16
%% paperName: 2020年山东等级考第 16 题 9 分
%% typeId: 6
%% type: 计算题
%% chapter: 曲线运动
%% point: 斜抛运动
%% method: 无
%% score: 9
%% degree: 450
%% duplicateId: 0
%% body:
单板滑雪 U 型池比赛是冬奥会比赛项目，其场地可以简化为如图甲所示的模型：U 形滑道由两个半径相同的四分之一圆柱面轨道和一个中央的平面直轨道连接而成，轨道倾角为 $17.2^\circ$。某次练习过程中，运动员以 $v_{M}=10\Ums$ 的速度从轨道边缘上的 M 点沿轨道的竖直切面 ABCD 滑出轨道，速度方向与轨道边缘线 AD 的夹角 $\alpha=72.8^\circ$，腾空后沿轨道边缘的 N 点进入轨道。图乙为腾空过程左视图。该运动员可视为质点，不计空气阻力，取重力加速度的大小 $g=10\Umsq$，$\sin72.8^\circ=0.96$，$\cos72.8^\circ=0.30$。求：
\begin{enumerate}
\item
运动员腾空过程中离开 AD 的距离的最大值 $d$；

\item
M、N 之间的距离 $L$。
\end{enumerate}

\twopicture[numA=甲, numB=乙, wA=6.00cm, wB=7.00cm]{figs/16a.png}{figs/16b.png}

%% answer:
\jdanswer{
\begin{enumerate}
\item
$4.8\Um$

\item
$12\Um$

\end{enumerate}
}
%% memo:
\memoanswer{
（1）在 M 点，设运动员在 ABCD 面内垂直 AD 方向的分速度为 $v_{1}$，由运动的合成与分解规律得 $v_{1}=v_{M}\sin72.8^\circ$；设运动员在 ABCD 面内垂直 AD 方向的分加速度为 $a_{1}$，由牛顿第二定律得 $mg\cos17.2^\circ=ma_{1}$；由运动学公式得 $d=\frac{v_{1}^{2}}{2a_{1}}$，联立并代入数据得 $d=4.8\Um$。

（2）在 M 点，设运动员在 ABCD 面内平行 AD 方向的分速度为 $v_{2}$，由运动的合成与分解规律得 $v_{2}=v_{M}\cos72.8^\circ$；设运动员在 ABCD 面内平行 AD 方向的分加速度为 $a_{2}$，由牛顿第二定律得 $mg\sin17.2^\circ=ma_{2}$；设腾空时间为 $t$，由运动学公式得 $t=\frac{2v_{1}}{a_{1}}$，$L=v_{2}t+\frac{1}{2}a_{2}t^{2}$，联立并代入数据得 $L=12\Um$。
}

\item
%% number: 17
%% paperName: 2020年山东等级考第 17 题 14 分
%% typeId: 6
%% type: 计算题
%% chapter: 磁场
%% point: 带电粒子在磁场中的运动
%% method: 无
%% score: 14
%% degree: 450
%% duplicateId: 0
%% body:
某型号质谱仪的工作原理如图甲所示。M、N 为竖直放置的两金属板，两板间电压为 $U$，Q 板为记录板，分界面 P 将 N、Q 间区域分为宽度均为 $d$ 的 I、Ⅱ两部分，M、N、P、Q 所在平面相互平行，a、b 为 M、N 上两正对的小孔。以 a、b 所在直线为 $z$ 轴，向右为正方向，取 $z$ 轴与 Q 板的交点 $O$ 为坐标原点，以平行于 Q 板水平向里为 $x$ 轴正方向，竖直向上为 $y$ 轴正方向，建立空间直角坐标系 $Oxyz$。区域 I、Ⅱ 内分别充满沿 $x$ 轴正方向的匀强磁场和匀强电场，磁感应强度大小、电场强度大小分别为 $B$ 和 $E$。一质量为 $m$，电荷量为 $+q$ 的粒子，从 a 孔飘入电场（初速度视为零），经 b 孔进入磁场，过 P 面上的 c 点（图中未画出）进入电场，最终打到记录板 Q 上。不计粒子重力。
\begin{enumerate}
\item
求粒子在磁场中做圆周运动的半径 $R$ 以及 c 点到 $z$ 轴的距离 $L$；

\item
求粒子打到记录板上位置的 $x$ 坐标；

\item
求粒子打到记录板上位置的 $y$ 坐标（用 $R$、$d$ 表示）；

\item
如图乙所示，在记录板上得到三个点 s$_{1}$、s$_{2}$、s$_{3}$，若这三个点是质子 \ce{^{1}_{1}H}、氚核 \ce{^{3}_{1}H}、氦核 \ce{^{4}_{2}He} 的位置，请写出这三个点分别对应哪个粒子（不考虑粒子间的相互作用，不要求写出推导过程）。
\end{enumerate}

\twopicture[numA=甲, numB=乙, wA=6.50cm, wB=4.20cm]{figs/17a.png}{figs/17b.png}

%% answer:
\jdanswer{
\begin{enumerate}
\item
$R=\frac{\sqrt{2mqU}}{qB}$，$L=\frac{\sqrt{2mqU}}{qB}-\sqrt{\frac{2mU}{qB^{2}}-d^{2}}$

\item
$x=\frac{md^{2}E}{4mU-2qd^{2}B^{2}}$

\item
$y=R-\sqrt{R^{2}-d^{2}}+\frac{d^{2}}{\sqrt{R^{2}-d^{2}}}$

\item
s$_{1}$、s$_{2}$、s$_{3}$ 分别对应氚核 \ce{^{3}_{1}H}、氦核 \ce{^{4}_{2}He}、质子 \ce{^{1}_{1}H} 的位置

\end{enumerate}
}
%% memo:
\memoanswer{
（1）设粒子经加速电场到 b 孔的速度大小为 $v$，粒子在区域 I 中做匀速圆周运动对应圆心角为 $\alpha$。在 M、N 两金属板间，由动能定理得 $qU=\frac{1}{2}mv^{2}$；在区域 I 中，粒子做匀速圆周运动，磁场力提供向心力，由牛顿第二定律得 $qvB=m\frac{v^{2}}{R}$，联立解得 $R=\frac{\sqrt{2mqU}}{qB}$；由几何关系得 $d^{2}+(R-L)^{2}=R^{2}$，$\cos\alpha=\frac{\sqrt{R^{2}-d^{2}}}{R}$，$\sin\alpha=\frac{d}{R}$，联立解得 $L=\frac{\sqrt{2mqU}}{qB}-\sqrt{\frac{2mU}{qB^{2}}-d^{2}}$。

（2）设区域Ⅱ中粒子沿 $z$ 轴方向的分速度为 $v_{z}$，沿 $x$ 轴正方向加速度大小为 $a$，位移大小为 $x$，运动时间为 $t$，由牛顿第二定律得 $qE=ma$；粒子在 $z$ 轴方向做匀速直线运动，由运动合成与分解的规律得 $v_{z}=v\cos\alpha$，$d=v_{z}t$；粒子在 $x$ 方向做初速度为零的匀加速直线运动，由运动学公式得 $x=\frac{1}{2}at^{2}$，联立解得 $x=\frac{md^{2}E}{4mU-2qd^{2}B^{2}}$。

（3）设粒子沿 $y$ 方向偏离 $z$ 轴的距离为 $y$，其中在区域Ⅱ中沿 $y$ 方向偏离的距离为 $y'$，由运动学公式得 $y'=vt\sin\alpha$，由题意得 $y=L+y'$，联立解得 $y=R-\sqrt{R^{2}-d^{2}}+\frac{d^{2}}{\sqrt{R^{2}-d^{2}}}$。

（4）根据（2）中结果 $x=\frac{md^{2}E}{4mU-2qd^{2}B^{2}}=\frac{d^{2}E}{4U-2d^{2}B^{2}\frac{q}{m}}$，可知 $\frac{q}{m}$ 越大，粒子打在记录板上的位置 $x$ 的坐标值越大，可得 s$_{1}$、s$_{2}$、s$_{3}$ 分别对应氚核 \ce{^{3}_{1}H}、氦核 \ce{^{4}_{2}He}、质子 \ce{^{1}_{1}H} 的位置。
}

\item
%% number: 18
%% paperName: 2020年山东等级考第 18 题 16 分
%% typeId: 6
%% type: 计算题
%% chapter: 运动和力
%% point: 动量和能量
%% method: 无
%% score: 16
%% degree: 350
%% duplicateId: 0
%% body:
如图所示，一倾角为 $\theta$ 的固定斜面的底端安装一弹性挡板，P、Q 两物块的质量分别为 $m$ 和 $4m$，Q 静止于斜面上 A 处。某时刻，P 以沿斜面向上的速度 $v_{0}$ 与 Q 发生弹性碰撞。Q 与斜面间的动摩擦因数等于 $\tan\theta$，设最大静摩擦力等于滑动摩擦力。P 与斜面间无摩擦，与挡板之间的碰撞无动能损失。两物块均可以看作质点，斜面足够长，Q 的速度减为零之前 P 不会与之发生碰撞。重力加速度大小为 $g$。
\begin{enumerate}
\item
求 P 与 Q 第一次碰撞后瞬间各自的速度大小 $v_{P1}$、$v_{Q1}$；

\item
求第 $n$ 次碰撞使物块 Q 上升的高度 $h_{n}$；

\item
求物块 Q 从 A 点上升的总高度 $H$；

\item
为保证在 Q 的速度减为零之前 P 不会与之发生碰撞，求 A 点与挡板之间的最小距离 $s$。
\end{enumerate}

\onepicture[w=4.50cm, h=2.35cm, align=right]{figs/18.png}

%% answer:
\jdanswer{
\begin{enumerate}
\item
$v_{P1}=\frac{3}{5}v_{0}$，$v_{Q1}=\frac{2}{5}v_{0}$

\item
$h_{n}=\left(\frac{7}{25}\right)^{n-1}\cdot\frac{v_{0}^{2}}{25g}$（$n=1$，2，3$\cdots\cdots$）

\item
$H=\frac{v_{0}^{2}}{18g}$

\item
$s=\frac{(8\sqrt{7}-13)v_{0}^{2}}{200g\sin\theta}$

\end{enumerate}
}
%% memo:
\memoanswer{
（1）P 与 Q 的第一次碰撞，取 P 的初速度方向为正方向，由动量守恒定律得 $mv_{0}=mv_{P1}+4mv_{Q1}$，由机械能守恒定律得 $\frac{1}{2}mv_{0}^{2}=\frac{1}{2}mv_{P1}^{2}+\frac{1}{2}\cdot4mv_{Q1}^{2}$，联立解得 $v_{P1}=-\frac{3}{5}v_{0}$，$v_{Q1}=\frac{2}{5}v_{0}$，故第一次碰撞后 P 的速度大小为 $\frac{3}{5}v_{0}$，Q 的速度大小为 $\frac{2}{5}v_{0}$。

（2）设第一次碰撞后 Q 上升的高度为 $h_{1}$，对 Q 由运动学公式得 $0-v_{Q1}^{2}=2\cdot(-2g\sin\theta)\cdot\frac{h_{1}}{\sin\theta}$，联立解得 $h_{1}=\frac{v_{0}^{2}}{25g}$。

设 P 运动至与 Q 刚要发生第二次碰撞前的位置时速度为 $v_{02}$，第一次碰后至第二次碰前，对 P 由动能定理得 $\frac{1}{2}mv_{02}^{2}-\frac{1}{2}mv_{P1}^{2}=-mgh_{1}$，联立解得 $v_{02}=\frac{\sqrt{7}}{5}v_{0}$。

P 与 Q 的第二次碰撞，设碰后 P 与 Q 的速度分别为 $v_{P2}$、$v_{Q2}$，由动量守恒定律得 $mv_{02}=mv_{P2}+4mv_{Q2}$，由机械能守恒定律得 $\frac{1}{2}mv_{02}^{2}=\frac{1}{2}mv_{P2}^{2}+\frac{1}{2}\cdot4mv_{Q2}^{2}$，联立解得 $v_{P2}=-\frac{3}{5}\times\frac{\sqrt{7}}{5}v_{0}$，$v_{Q2}=\frac{2}{5}\times\frac{\sqrt{7}}{5}v_{0}$。

设第二次碰撞后 Q 上升的高度为 $h_{2}$，对 Q 由运动学公式得 $0-v_{Q2}^{2}=2\cdot(-2g\sin\theta)\cdot\frac{h_{2}}{\sin\theta}$，联立解得 $h_{2}=\frac{7}{25}\cdot\frac{v_{0}^{2}}{25g}$。

设 P 运动至与 Q 刚要发生第三次碰撞前的位置时速度为 $v_{03}$，第二次碰后至第三次碰前，对 P 由动能定理得 $\frac{1}{2}mv_{03}^{2}-\frac{1}{2}mv_{P2}^{2}=-mgh_{2}$，联立解得 $v_{03}=\left(\frac{\sqrt{7}}{5}\right)^{2}v_{0}$。

P 与 Q 的第三次碰撞，设碰后 P 与 Q 的速度分别为 $v_{P3}$、$v_{Q3}$，由动量守恒定律得 $mv_{03}=mv_{P3}+4mv_{Q3}$，由机械能守恒定律得 $\frac{1}{2}mv_{03}^{2}=\frac{1}{2}mv_{P3}^{2}+\frac{1}{2}\cdot4mv_{Q3}^{2}$，联立解得 $v_{P3}=-\frac{3}{5}\times\left(\frac{\sqrt{7}}{5}\right)^{2}v_{0}$，$v_{Q3}=\frac{2}{5}\times\left(\frac{\sqrt{7}}{5}\right)^{2}v_{0}$。

设第三次碰撞后 Q 上升的高度为 $h_{3}$，由运动学公式得 $0-v_{Q3}^{2}=2\cdot(-2g\sin\theta)\cdot\frac{h_{3}}{\sin\theta}$，联立解得 $h_{3}=\left(\frac{7}{25}\right)^{2}\cdot\frac{v_{0}^{2}}{25g}$。

总结可知，第 $n$ 次碰撞后，物块 Q 上升的高度为 $h_{n}=\left(\frac{7}{25}\right)^{n-1}\cdot\frac{v_{0}^{2}}{25g}$（$n=1$，2，3$\cdots\cdots$）。

（3）当 P、Q 达到 $H$ 时，两物块到此处的速度可视为零，对两物块运动全过程由动能定理得 $0-\frac{1}{2}mv_{0}^{2}=-(m+4m)gH-\tan\theta\cdot4mg\cos\theta\cdot\frac{H}{\sin\theta}$，解得 $H=\frac{v_{0}^{2}}{18g}$。

（4）设 Q 第一次碰撞至速度减为零需要的时间为 $t_{1}$，由运动学公式得 $v_{Q1}=2gt_{1}\sin\theta$；设 P 运动到斜面底端时的速度为 $v_{P1}'$，需要的时间为 $t_{2}$，由运动学公式得 $v_{P1}'=v_{P1}+gt_{2}\sin\theta$，$v_{P1}'^{2}-v_{P1}^{2}=2sg\sin\theta$；设 P 从 A 点到 Q 第一次碰后速度减为零处匀减速运动的时间为 $t_{3}$，得 $v_{02}=(-v_{P1})-gt_{3}\sin\theta$；当 A 点与挡板之间的距离最小时 $t_{1}=2t_{2}+t_{3}$，联立并代入数据得 $s=\frac{(8\sqrt{7}-13)v_{0}^{2}}{200g\sin\theta}$。
}

\end{enumerate}

\end{document}
